Long-context language models are increasingly used to summarize records that are too long for routine human inspection. For these settings, the relevant question is not only whether summaries are faithful, but which error types grow as source length and fact load increase. We study this question with controlled long-document summarization tasks built from \qmsum and \govreport background text. We inject known compliance-incident facts as ground truth, then ask \gptfourmini and \gptfivemini to return structured summaries across 48 main runs spanning roughly 2k, 6k, 12k, and 24k source tokens. The experiment includes fixed-load tasks with 10 target facts at every length and scaled-load tasks with 4, 8, 16, and 32 target facts. Across 600 target records and 536 returned records, hallucinations and duplicates are zero, normalized unit, object, count, month, and severity errors are zero, and only two action-field errors occur. The only strong length-dependent failure is structural: \gptfivemini returns empty visible content in two 24k-token, 32-fact runs under a 5,000 completion-token cap. Rerunning those cases with a 12,000-token cap recovers all records with zero omissions or field errors. These results suggest that, in this controlled setting, long-context summarization fails abruptly at a task-budget boundary rather than through gradual factual drift.
